\begin{pastexam}[label={px:ncu-109-dsa-11}]{中央資工所 109 學年度 第 11 題}
（是非；答對 1 分，答錯倒扣 0.5 分）

If $V(G')\subseteq V(G)$ and $E(G')\subseteq E(G)$, a graph G' is G's connected component.\\
(A)~True. \qquad (B)~False.
\end{pastexam}

\begin{pastsoln}
答案：(B) False。

$V(G') \subseteq V(G)$ 而且 $E(G') \subseteq E(G)$ 只說明 $G'$ 是 $G$ 的\textbf{子圖}（subgraph）。
\textbf{連通元件}（connected component）要求更多（資結書第~21.1~節）：它是 $G$ 的\textbf{極大}連通
子圖，也就是 $G'$ 本身連通，而且再加進 $G$ 的任何一個頂點或邊都不再連通（或者說，和 $G'$ 有路相通的
頂點、邊全都已經在 $G'$ 裡）。

反例：$G$ 是一條路徑 $1$--$2$--$3$。 取 $G'$ 只有頂點 $\{1, 2\}$、邊 $\{(1, 2)\}$，兩個包含關係都成立，
$G'$ 也連通，但 $3$ 和 $2$ 相連卻不在 $G'$ 裡，不是極大的；$G$ 唯一的連通元件是整條路徑。 再取 $G'$
為頂點 $\{1, 3\}$、沒有邊，它本身就不連通，更不是連通元件。
\end{pastsoln}
